第一百四十一类是立场先于检验段：先写决策者更在意哪类错误，再确定原假设、备择假设、功效
和样本量。第一百四十二类是无样本先交代模拟器段：公开缺失的数据、模拟分布和用途，只把
模拟结果限定为流程比较。第一百四十三类是结构规模授权搜索段：列出节点、边或方案数量，
说明精确枚举负担后再切换搜索方法。第一百四十四类是下游后果触发预处理段：不先说“去噪”，
而先说明干扰会怎样改变反演或决策，再给处理规则。第一百四十五类是有限改善收窄结论段：
主量基本不变而局部指标改善时，只称局部修正，不写“全面提升”。

第一百四十六类是共同环境复评段：各候选先在各自情景中生成，再放入同一随机环境比较，避免
把情景差异误当方案差异。第一百四十七类是管理语言参数化段：把“优先、稳定、兼顾”等口号
翻译为可调参数或优先级，并说明参数改变哪条约束。第一百四十八类是安全--及时--可执行损失段：
先分别定义三种错误后果，再选择阈值或动态规划目标。第一百四十九类是样本行转主体段：重复
测量资料先聚合到真实决策主体，再计算首次达标时点，避免行数替代人数。第一百五十类是总体--
少数类--样本数并列段：分类结果同时给总体指标、关键少数类召回和对应样本量，不用单一准确率
结束段落。

句子节奏也有可辨认的规律。定义和假设用短句，避免一个句子同时承担五个动作；
推导段允许稍长，但每个公式后立即解释变量或几何意义；结果段先给条件和数字，
下一句再说明趋势或机制。图表不以“如下图所示”结束，而写成“由图可知，温度升高
后……先……后……”，或按“图左/图右”分别完成比较。相比之下，“构建综合性框架”
“充分证明模型有效”“显著提升了性能”“为后续研究奠定基础”等 AI 惯用语若没有
具体对象、数值和证据，一律不用。

上述结论由 59/59 篇论文的文风入口支持：其中四十四篇保留完成原页核验的细粒度
句式和段落节奏，十五篇补入结构化判断动作卡。它们的作用是提供可调用的写作动作，
不把人工选出的高信息段落冒充 59 篇全文词频。每条惯用表达均与观察、限制、动作或结果绑定；
若省略处填不进本题事实，便不调用该句式。

